\documentclass[11pt,a4paper]{article}

%% PACHETE MATEMATICE
\usepackage{amsmath,amssymb,amsfonts}
\usepackage{amsthm}
\usepackage{mathpazo}
\usepackage{eulervm}
\usepackage{enumerate}
\usepackage{color}
\usepackage{graphicx}
\usepackage[all]{xy}
\usepackage{xcolor}
\usepackage{framed}
\usepackage[unicode]{hyperref}
\setcounter{MaxMatrixCols}{20}
\usepackage{multicol}
%\usepackage[romanian]{babel}


%% ASPECT
\linespread{1.05}
\usepackage[T1]{fontenc}
\usepackage[utf8x]{inputenc}
\usepackage[margin=2cm]{geometry}
% \usepackage[color]{showkeys}
\usepackage{anyfontsize}
\usepackage{titlesec}
\titleformat*{\section}{\large\bfseries}

\usepackage{fancyhdr}
\pagestyle{fancy}
\rhead{\small \color{gray} Notițe de Adrian Manea}
\lhead{\small \color{gray} ETTI-AG, 2017---2018, A. Niță \& A. Oprina}


\usepackage[autostyle = false, style = german]{csquotes}
\usepackage[romanian]{babel}
\newcommand\qq{\enquote}

%% COMENZILE MELE
\renewcommand*{\proofname}{Dem.:}
\newcommand\bb{\mathbb}
\newcommand\dr{\mathrm}
\newcommand\kal{\mathcal}
\newcommand\fk{\mathfrak}
\newcommand\op{\oplus}
\newcommand\ot{\otimes}
\newcommand\ds{\displaystyle}
\newcommand\id{\indent}
\newcommand\nid{\noindent}
\newcommand\seq{\subseteq}
\newcommand\sneq{\subsetneq}
\newcommand\speq{\supseteq}
\newcommand\spneq{\supsetneq}
\newcommand\rar{\rightarrow}
\newcommand\Rar{\Rightarrow}
\newcommand\xrar{\xrightarrow}
\newcommand\lar{\leftarrow}
\newcommand\Lar{\Leftarrow}
\newcommand\xlar{\xleftarrow}
\newcommand\lrar{\leftrightarrow}
\newcommand\Lrar{\Leftrightarrow}
\newcommand\ideal{\trianglelefteq}
\newcommand\ai{\text{ a.\^{i}. }}
\newcommand\sk{(\textit{Schi\c{t}\u{a}}) }
\newcommand\rhup{\rightharpoonup}
\newcommand\rhdn{\rightharpoondown}
\newcommand\lhup{\leftharpoonup}
\newcommand\lhdn{\leftharpoondown}
\newcommand\vid{\emptyset}
\newcommand\wh{\widehat}
\newcommand\ol{\overline}
\newcommand\NN{\mathbb{N}}
\newcommand\RR{\mathbb{R}}
\newcommand\ZZ{\mathbb{Z}}
\newcommand\QQ{\mathbb{Q}}
\newcommand\CC{\mathbb{C}}

%% STILURILE MELE
\newtheoremstyle{dotless}{}{}{\itshape}{}{\bfseries}{:}{ }{}

\theoremstyle{dotless}
\newtheorem{theorem}{Teorem\u{a}}[section]
\newtheorem{proposition}{Propozi\c{t}ie}[section]
\newtheorem{lemma}{Lem\u{a}}[section]
\newtheorem{corollary}{Corolar}[section]
\newtheorem{exercise}{Exerci\c{t}iu}

\newtheoremstyle{dotlessdr}{}{}{}{}{\bfseries}{:}{ }{}
\theoremstyle{dotlessdr}
\newtheorem{example}{Exemplu}[section]
\newtheorem{definition}{Defini\c{t}ie}[section]
\newtheorem{remark}{Observa\c{t}ie}[section]




\begin{document}

\begin{center}
  {\large\textbf{Seminar 5 \\ Vectori și valori proprii}}
\end{center}

\vspace{1cm}

\section{Vectori și valori proprii}
\label{sec:vvp}

Fie $V$ un spațiu vectorial peste un corp $\Bbbk$ și $f : V \to V$ o aplicație
liniară.

\begin{definition}\label{def:vvp}
  Un vector $v \in V$ se numește \emph{vector propriu} (eng. \emph{eigenvector})
  pentru aplicația $f$ dacă există un scalar $\lambda \in \Bbbk$ astfel încît
  $f(v) = \lambda \cdot v$.

  În acest caz, $\lambda$ se numește \emph{valoarea proprie} (eng. \emph{eigenvalue})
  asociată vectorului  propriu $v$.
\end{definition}

Așadar, vectorii proprii sînt aceia pentru care aplicația are o acțiune simplă,
de \enquote{rescalare}, adică doar înmulțire cu un scalar, care se numește
valoarea proprie asociată.

În exerciții, pentru a găsi valorile proprii asociate unei aplicații liniare
ca mai sus (endomorfism), se procedează astfel:
\begin{itemize}
\item Presupunem că $\dim_\Bbbk V = n$. Scriem matricea aplicației $f$ în
  baza canonică a lui $V$, $M^B(f) \in M_n(\Bbbk)$;
\item Scriem \emph{polinomul caracteristic} al matricei, anume:
  \[
    P(x) = \det(A - x \cdot I_n).
  \]
\item Rădăcinile polinomului caracteristic se vor nota cu $\lambda_i$ și
  sînt valorile proprii ale endomorfismului.
\end{itemize}

Mulțimea valorilor proprii se mai numește \emph{spectrul} endomorfismului
$f$, notat uneori cu $\sigma(f)$.

Apoi, pentru a găsi vectorii proprii asociați valorilor proprii de mai sus
se rezolvă, pentru fiecare $\lambda_i$, ecuația $f(v) = \lambda_i v$ și
se găsește $v$.

Cu ajutorul acestor vectorii proprii, obținem următoarele:

\begin{definition}\label{def:sp-invar}
  Fie $\lambda$ o valoare proprie a unui endomorfism $f$.

  Dacă $v$ este un vector propriu asociat valorii proprii $\lambda$,
  $V(\lambda) = V_\lambda = \mathrm{Sp}(v)$ este un subspațiu al lui
  $V$, numit \emph{subspațiul invariant} (propriu) generat de $v$.
\end{definition}

Pentru o valoare proprie $\lambda$, se definesc:

\begin{definition}\label{def:ma-mg}
  Fie $\lambda$ o valoare proprie a endomorfismului $f$. Se numește
  \emph{multiplicitatea algebrică} a lui $\lambda$, notată $m_a(\lambda)$
  sau $a_\lambda$ puterea la care apare factorul $(X - \lambda)$ în descompunerea
  polinomului caracteristic al lui $f$.

  Se numește \emph{multiplicitatea geometrică} a lui $\lambda$, notată $m_g(\lambda)$
  sau $g_\lambda$ dimensiunea subspațiului $V_\lambda$.
\end{definition}

Au loc următoarele proprietăți utile:
\begin{proposition}\label{pr:vvp}
  Fie $A$ matricea endomorfismului $f$ și fie $\lambda_i$ valorile sale proprii.

  \begin{itemize}
  \item Suma valorilor proprii (cu tot cu multiplicități) este egală cu urma
    matricei, iar produsul lor este egal cu determinantul matricei.

    În particular, dacă $0$ este valoare proprie, atunci $A$ este neinversabilă;
  \item $\lambda_i^k$ este valoare proprie pentru $A^k$, pentru orice $k \geq 1$;
  \item \textbf{(Teorema Cayley-Hamilton)} $P_A(A) = 0$, unde $P_A$ este polinomul
    caracteristic al matricei $A$.
  \end{itemize}
\end{proposition}

%%%%%%%%%%%%%%%%%%%%%%%%%%%%%%%%%%%%%%%%%%%%%%%%%%%%%%%%%%%%%%%%%%%%%%

\section{Diagonalizare}
\label{sec:diag}

\begin{definition}\label{def:diag}
  O matrice $A \in M_n(\Bbbk)$ se numește \emph{diagonalizabilă} dacă ea este
  asemenea cu o matrice diagonală $D$, i.e.\ care are elemente nenule doar pe
  diagonala principală.

  Altfel spus, există $T \in M_n(\Bbbk)$ inversabilă, astfel încît $T^{-1}AT = D$.
\end{definition}

\begin{remark}\label{rk:diag-vvp}
  Dacă $A$ este diagonalizabilă, atunci $D$ (din definiția de mai sus) conține pe
  diagonală doar valorile proprii ale lui $A$.
\end{remark}

Pașii pentru verificarea și obținerea diagonalizării unei matrice sînt:
\begin{enumerate}[(1)]
\item Fixăm o bază $B$ oarecare a lui $V$ și scriem matricea lui $f$ în
  baza $B$;
\item Determinăm valorile proprii $\lambda_1, \dots, \lambda_p$, cu ordinele
  de multiplicitate $n_1, \dots, n_p$;
\item Pentru fiecare valoare proprie $\lambda_i$, se determină subspațiul propriu
  (invariant) $V_{\lambda_i}$ și o bază a sa, $B_i$;
\item Dacă există o valoare proprie $\lambda_i$, astfel încît
  $\dim V_{\lambda_i} \neq n_i$, algoritmul se oprește, deoarece matricea nu
  se poate diagonaliza;
\item În caz contrar, dacă avem $a_{\lambda_i} = g_{\lambda_i}$, atunci
  matricea se poate diagonaliza, iar $B' = B_1 \cup \dots \cup B_p$ este
  o bază a spațiului $V$, formată din vectori proprii;
\item Se determină matricea $T$, de trecere de la baza $B$ la baza $B'$, care
  este inversabilă. În plus, obținem:
  \[
    M^{B'}(f) = T^{-1} M^B(f) T = \mathrm{diag}(\lambda_1, \dots, \lambda_p),
  \]
  adică o matrice diagonală cu valorile proprii $\lambda_i$ pe diagonală.
\end{enumerate}


%%%%%%%%%%%%%%%%%%%%%%%%%%%%%%%%%%%%%%%%%%%%%%%%%%%%%%%%%%%%%%%%%%%%%%

\section{Exerciții}
\label{sec:exc}

1. Să se determine vectorii și valorile proprii pentru matricele:
\[
  A = \begin{pmatrix}
    1 & 0 \\
    8 & 2
  \end{pmatrix}, \quad
  B = \begin{pmatrix}
    0 & 2 & 2 \\
    2 & 0 & -2 \\
    2 & -2 & 0
  \end{pmatrix}, \quad
  C = \begin{pmatrix}
    0 & 1 & -2 \\
    1 & 1 & 1 \\
    1 & -1 & -1
  \end{pmatrix}.
\]

\vspace{1cm}

2. Fie aplicația liniară $f : \RR^3 \to \RR^3, f(x, y, z) = (-y, x, z)$.

Să se calculeze vectorii și valorile proprii ale lui $f$.

\vspace{1cm}

3. Fie aplicația liniară:
\[
  f : \RR^3 \to \RR^3, \quad f(x, y, z) = (x - z, 8x + y - 2z, z).
\]
Să se calculeze vectorii și valorile proprii ale lui $f$.

\vspace{1cm}

4. Fie aplicația liniară $f : M_2(\RR) \to M_2(\RR)$, definită prin:
\[
  f \begin{pmatrix}
    a & b \\
    c & d
  \end{pmatrix} = \begin{pmatrix}
    2a + b + d & 2b + 4c + 5d \\
    2c + d & 8d
  \end{pmatrix}
\]

Să se arate că $f$ nu este diagonalizabilă.

\vspace{1cm}

5. Să se diagonalizeze endomorfismul $f : \RR^3 \to \RR^3$, definit prin:
\[
  f(x, y, z) = (y + z, x + z, x + y).
\]

\vspace{1cm}

6. Să se diagonalizeze matricea:
\[
  A = \begin{pmatrix}
    7 & 2 & -2 \\
    2 & 4 & -1 \\
    -2 & -1 & 4
  \end{pmatrix}
\]
și apoi, folosind forma diagonală, să se calculeze $A^n$, cu $n \in \NN$.

\vspace{1cm}

7. Aceeași cerință pentru matricele:
\[
  B = \begin{pmatrix}
    1 & 1 \\
    4 & 1 
  \end{pmatrix}, \quad
  C = \begin{pmatrix}
    11 & -5 & 5 \\
    -5 & 3 & -3 \\
    5 & -3 & 3
  \end{pmatrix}.
\]

\vspace{1cm}

8. Determinați vectorii și valorile proprii pentru aplicațiile:
\begin{enumerate}[(a)]
\item $f : \RR[X]_3 \to \RR[X]_3, f(P) = 2P'$;
\item $f : \RR^3 \to \RR^3, f(a, b, c) = (a - b + c, b, b - c)$.
\end{enumerate}

\vspace{1cm}

9. Să se diagonalizeze endomorfismul $f : \RR^4 \to \RR^4$, definit prin:
\[
  f(x, y, z, t) = (x + y + z + t, x + y - z - t, x - y + z - t, x - y - z + t).
\]
\end{document}
